\section{Termékcélok}

A termékcélok meghatározása a PRD egyik kulcsfontosságú része, mert ezek jelölik ki, hogy a fejlesztés során milyen értéket kell létrehozni a felhasználók számára. A célok nem azonosak a konkrét funkciókkal: inkább azt fogalmazzák meg, hogy az alkalmazásnak milyen problémára kell megoldást adnia, milyen felhasználói élményt kell támogatnia, és milyen szempontok alapján lehet eldönteni, hogy az első verzió sikeresnek tekinthető-e.

A WatchDB elsődleges termékcélja, hogy a felhasználó gyorsan és egyszerűen tudjon filmeket keresni, majd azokat saját listáihoz hozzáadni. Ez a cél közvetlenül kapcsolódik a projekt kiinduló problémájához: a filmnézési élmények több platformon és helyzetben történnek, ezért a felhasználónak szüksége van egy egységes, platformfüggetlen nyilvántartásra. Az MVP-ben ennek a célnak a filmkeresés, a filmrészletek megjelenítése, a watchlist és a megnézett filmek naplózása felel meg.

A második fontos cél a személyes filmnapló felépítése. A WatchDB nem csupán egy megnézendő filmeket tartalmazó lista, hanem olyan személyes filmnapló, amelyben a felhasználó visszakeresheti, milyen filmeket látott, hogyan értékelte őket, és milyen filmeket szeretne később megnézni. Ez a cél azért fontos, mert az alkalmazásnak akkor is értéket kell adnia, ha a felhasználó egyelőre nem használ közösségi funkciókat. Az MVP tehát önmagában is használható termékként jelenik meg, nem csak egy későbbi, nagyobb rendszer előzetes változataként.

A harmadik termékcél a közösségi filmfelfedezés előkészítése. A teljes termékvízió alapján a WatchDB hosszabb távon lehetőséget adna arra, hogy a felhasználók kövessék egymás aktivitását, megtekintsék barátaik értékeléseit, és ezek alapján fedezzenek fel új filmeket. Ez a cél az MVP-ben még nem feltétlenül jelenik meg teljes funkcionalitásként, de fontos tervezési szempontként hat a rendszerre. A felhasználói profil, a naplózott filmek strukturált tárolása és a később bővíthető adatmodell mind olyan elemek, amelyek támogatják ezt az irányt.

A negyedik cél a gyors, reszponzív és könnyen használható felhasználói élmény biztosítása. Mivel a célközönség jelentős része alkalmi filmnéző, az alkalmazásnak nem szabad túl sok döntést vagy bonyolult lépést kérnie a felhasználótól. A film keresése, listához adása vagy értékelése akkor működik jól, ha néhány kattintással elvégezhető. Ez a cél nemcsak felületi kérdés, hanem technikai döntéseket is befolyásol: fontos a gyors oldalbetöltés, a hatékony keresés, a mobilon is használható felület és az egyértelmű navigáció.

A PRD-ben szereplő sikerességi mutatók ezekhez a célokhoz kapcsolódnak. Ilyen mutató lehet például a regisztrált felhasználók száma, az aktív felhasználók aránya, az egy felhasználóra jutó havi naplózott filmek száma, a kereséstől a naplózásig eltelt idő, illetve későbbi verzióban az is, hogy hány felhasználó kapcsolódik legalább egy ismerőshöz. Ezek a mérőszámok azért hasznosak, mert a termékcélokat ellenőrizhetővé teszik: nemcsak azt mutatják meg, hogy elkészültek-e a funkciók, hanem azt is, hogy a felhasználók valóban használják-e őket. Emellett a termékcélok a későbbi iteratív fejlesztésben is támpontot adhatnak: segíthetnek eldönteni, hogy érdemes-e új funkcióval bővíteni az alkalmazást, vagy éppen azt, hogy egy kevéssé használt funkció fenntartása továbbra is indokolt-e.

Összességében a WatchDB termékcéljai a személyes értékteremtésre és a későbbi közösségi bővíthetőségre épülnek. Az első verzióban a hangsúly a gyors filmkeresésen, a watchlist kezelésén, a megnézett filmek naplózásán és az egyszerű értékelésen van. Ezek a célok elég szűkek ahhoz, hogy egy MVP keretében megvalósíthatók legyenek, ugyanakkor elég stabil állapotot adnak ahhoz, hogy a termék később közösségi filmfelfedező platformmá fejlődhessen.
